要证 \(\frac1n\to0\)，做法是等对方报出误差容限，再回一个序号：对 \(\varepsilon=0.01\) 取 \(N=101\)，此后所有项都落在带子里。量词 \(\forall\varepsilon>0\ \exists N\ \forall n\ge N\) 规定的正是这个先后次序——谁先出题，谁后应答。本章两篇都围着这个次序打转。

函数列把问题提高一档。此时每一项是一条曲线，误差要在定义域的每个点上分别检验，于是 \(N\) 多了一个可以依赖的对象：点 \(x\)。允许它依赖，得到点态收敛；不允许，得到一致收敛。两个定义的符号一模一样，只是 \(\exists N\) 与 \(\forall x\) 换了位置，而这一换决定了极限函数还是不是连续的、极限能不能与积分交换。

一个算得出数字的例子把差距摆明白。在 \([0,1]\) 上取 \(f_n(x)=x^n\)，要让误差小于 \(0.1\)：\(x=0.99\) 时需要 \(n\ge230\)，\(x=0.999\) 时需要 \(n\ge2302\)，\(x\) 越靠近 1 门槛越高且没有上限。每个点最终都达标，却不存在一个对所有点同时达标的时刻。逐点的好消息与整体的坏消息可以同时为真。

\subsection{怎么读这一章}

第一篇把``最终任意接近''翻译成量词，并把 \(\varepsilon\)--\(N\) 证明的标准动作练熟：先估误差、再反解 \(N\)。第二篇只改动一个量词的位置，就分出两种强度。读的时候反复问同一句话：控制误差的那个阶段，能不能对所有输入统一选定？许多性质能否传给极限，答案都压在这句话上。

\subsection{本章篇目}

\begin{itemize}
\tightlist
\item \hyperref[sec:12-Real-Analysis-Support/02-Convergence/01]{``用量词精确刻画数列趋近''}：定义、标准证法，以及唯一性、有界性、子列这些直接推论；
\item \hyperref[sec:12-Real-Analysis-Support/02-Convergence/02]{``点态与一致收敛只差一个量词顺序''}：\(x^n\) 的边界层、sup 范数的统一发言，以及一致收敛管得到与管不到的地方。
\end{itemize}

带着这个问题往下走，后面几章会不断改写它的形式：把定义域换成一般度量空间，把 sup 换成积分意义下的距离，把``对所有点统一''换成``除掉一个零测集后统一''。次序之争的形态在变，争的东西没变。
